\chapter{The Dihedral Polygon Census}
\label{chap:harmonic-taxonomy}

\section{Census from \texorpdfstring{$\mathbb{Z}_1$ to $\mathbb{Z}_{16}$}{Z1 to Z16}}

Table~\ref{tab:poly-census-summary} details the complete census of
non-equivalent dihedral and cyclic inscribed polygon orbits across the
circular rings $\mathbb{Z}_1$ through $\mathbb{Z}_{16}$.

% Summary Table of Inscribed Polygon Orbits Z_1 to Z_16
\begin{table}[htbp]
\centering\small
\caption{\label{tab:poly-census-summary}Dihedral Polygon Orbit Census from $\mathbb{Z}_1$ to $\mathbb{Z}_{16}$}
\begin{tabular}{r | r r | r r | r r}
\toprule
Ring $\mathbb{Z}_n$ & $2^n$ Subsets & $|D_n|$ & Dihedral Orbits & Cyclic Orbits & Chiral Pairs & Chiral \% \\
\midrule
$\mathbb{Z}_{1}$ & 2 & 2 & \textbf{2} & 2 & 0 & 0.0\% \\
$\mathbb{Z}_{2}$ & 4 & 4 & \textbf{3} & 3 & 0 & 0.0\% \\
$\mathbb{Z}_{3}$ & 8 & 6 & \textbf{4} & 4 & 0 & 0.0\% \\
$\mathbb{Z}_{4}$ & 16 & 8 & \textbf{6} & 6 & 0 & 0.0\% \\
$\mathbb{Z}_{5}$ & 32 & 10 & \textbf{8} & 8 & 0 & 0.0\% \\
$\mathbb{Z}_{6}$ & 64 & 12 & \textbf{13} & 14 & 1 & 7.7\% \\
$\mathbb{Z}_{7}$ & 128 & 14 & \textbf{18} & 20 & 2 & 11.1\% \\
$\mathbb{Z}_{8}$ & 256 & 16 & \textbf{30} & 36 & 6 & 20.0\% \\
$\mathbb{Z}_{9}$ & 512 & 18 & \textbf{46} & 60 & 14 & 30.4\% \\
$\mathbb{Z}_{10}$ & 1,024 & 20 & \textbf{78} & 108 & 30 & 38.5\% \\
$\mathbb{Z}_{11}$ & 2,048 & 22 & \textbf{126} & 188 & 62 & 49.2\% \\
$\mathbb{Z}_{12}$ & 4,096 & 24 & \textbf{224} & 352 & 128 & 57.1\% \\
$\mathbb{Z}_{13}$ & 8,192 & 26 & \textbf{380} & 632 & 252 & 66.3\% \\
$\mathbb{Z}_{14}$ & 16,384 & 28 & \textbf{687} & 1,182 & 495 & 72.1\% \\
$\mathbb{Z}_{15}$ & 32,768 & 30 & \textbf{1224} & 2,192 & 968 & 79.1\% \\
$\mathbb{Z}_{16}$ & 65,536 & 32 & \textbf{2250} & 4,116 & 1,866 & 82.9\% \\
\bottomrule
\end{tabular}
\end{table}


The census reveals a phase transition in the emergence of chiral asymmetry:
\begin{itemize}
    \item On small rings $\mathbb{Z}_1$ through $\mathbb{Z}_5$, every
    inscribed polygon is strictly bilaterally symmetric ($0\%$ chirality).
    \item Chirality first emerges on $\mathbb{Z}_6$ at size $k=3$:
    the unique scalene triangle $\texttt{/013/}$, whose reflection
    $\texttt{/023/}$ forms an enantiomorphic partner.
    \item Chirality expands steadily through $\mathbb{Z}_8$ ($20.0\%$),
    crosses the $50\%$ majority mark at $\mathbb{Z}_{12}$ ($57.1\%$),
    and reaches $82.9\%$ at $\mathbb{Z}_{16}$, where $1{,}866$ of the
    $2{,}250$ orbits are chiral enantiomorphs.
\end{itemize}

\section{Analytical Invariants and Digital Resource}

Each orbit is uniquely identified by its canonical circular stride
composition word $c_1 \dots c_k$ (with digit sum $n$ and length
$k$), local atlas designation $k\text{-}i$, canonical vertex subset
representative in hexadecimal slash notation $/v_1 \dots v_k/$, dihedral
orbit length $|\mathcal{O}|$, and binary bilateral reflection flag ($1 =
\text{achiral}$, $0 = \text{chiral}$).

The complete 5{,}099-orbit inventory is maintained in the companion
files \path{data/inngons.tsv}, \path{inngons.csv}, and \path{inngons.json}
for computational querying and stream filtering.
